\section{Rules of evolution}\label{sec:laws}

The sequence of Section~\ref{sec:history} is not arbitrary. Reading the twelve steps as graph operations, the same few rules fire again and again. We state each rule in graph terms, give the step(s) in which it fires, and say what it would predict for an instance not yet built. Notation: $N^-_M(o)$ is the set of model nodes with an edge into the observable $o$ (framework assignments and candidate edges excluded); $\dim[g]$ is the mass dimension of the coupling of a model node; a \emph{route} $(L,m)$ has legs $L\subset V_S\times\mathbb N$ and head $m\in V_M$. The rules are regularities of this instance and are written so that they can fail on the next.

\paragraph{Rule 1 (symmetry partition).} \emph{If a framework node represents a symmetry group $G$ unitarily on a state node, the state node is replaced by the irreducible components of the representation, and the labels of the components (the Casimir invariants) become observables.} Fires at Step~2: $U(\Lambda,a)$ partitions $\mathcal H_1$ into Wigner's classes; $P^2=m^2$ and $W^2=-m^2s(s+1)$ become observables with badge P. The irreducible components are the \emph{building blocks} of the state layer: no model node maps one of them into a proper subspace of itself. Prediction for other instances: wherever a symmetry is imposed on the state layer (isospin SU(2), flavour SU(3), colour SU(3)), the state nodes split into multiplets and a new label becomes an observable.

\paragraph{Rule 2 (in-degree-zero drive).} \emph{A model node is added only when there exists an observable node $o$ with $N^-_M(o)=\emptyset$, or when the existing model layer is inconsistent on theory-internal grounds. In the first case the new node must send an edge to $o$; in the second it must reproduce every edge of the node it replaces.} Fires at Steps~3$\to$4 (particle production $\to$ field node), 5 (Compton, $a_e$, Lamb $\to$ QED vertex), 6 ($\beta$ spectrum $\to$ Fermi), 7 (muon lifetime $\to$ second contact node), 8 (parity asymmetries $\to$ V$-$A), 11 (oscillation $\to$ candidate node); the theory-internal case fires at Steps~9 (non-renormalizability and unitarity $\to$ gauge vertices) and 10 (forbidden mass terms $\to$ Higgs sector). In the graph, the red dashed ring is a to-do item for the model layer. The rule also says what \emph{never} happens: no model node is added with no observable to reach and no inconsistency to cure.

\paragraph{Rule 3 (Fock extension).} \emph{If the repairing model node must connect states of different particle number, the state layer is extended by the model node's back-edge: every irreducible block $\mathcal H_s$ becomes the ladder $\bigoplus_N \mathrm{Sym}^N\mathcal H_s$ or $\bigoplus_N\Lambda^N\mathcal H_s$, and no new irreducible block is introduced.} Fires at Step~4. Two features are worth stating. The extension is \emph{minimal}: it composes existing building blocks (tensor powers) rather than inventing new ones, which is why the sectors of Step~4 are exactly the representations of Step~2 with a ladder attached. And the trigger is \emph{relativistic}: the three observations of Step~3 involve energies at or above a creation threshold ($E\ge mc^2$ for pairs, the mass difference of the nucleus in $\beta$ decay, photon energy in emission). Non-relativistic observables never force this extension, which is why one-particle quantum mechanics was complete for atomic spectroscopy. Prediction: any instance whose observables involve number-changing transitions (phonons in solids, photons in cavity QED) must acquire ladders by the same operation.

\paragraph{Rule 4 (route composition).} \emph{Every interaction node $m$ is entered by routes whose legs are the (sector, occupation-number) pairs of the fields in the Lagrangian term of $m$; the occupation numbers are the minimal changes the term can effect; and the number of distinct routes into $m$ equals the number of sector combinations the term couples.} Fires from Step~5 on. The QED vertex has one route ($e(1)\wedge\gamma(1)$); Fermi's node one four-leg route; the V$-$A node three routes (the three processes it unifies); the charged-current vertex three routes (two lepton doublets and one quark doublet); the neutral-current vertex one route with a group leg (all fermions). The rule turns a Lagrangian into a set of hyperedges mechanically, and it is the rule an AI system can apply without judgement: the legs are read off the field content of the term.

\paragraph{Rule 5 (effective status).} \emph{A model node whose coupling has $\dim[g]<0$ is effective: it has a validity scale $\Lambda\sim|g|^{-1/|\dim g|}$, its observable edges are valid only for $q^2\ll\Lambda^2$, and it cannot be the final form of the model layer.} Fires at Steps~6--8 ($G_F$, $\dim=-2$, $\Lambda$ of a few hundred GeV). The rule is the graph form of the statement that non-renormalizable terms are low-energy limits; it marks the dashed nodes as temporary the moment they are added.

\paragraph{Rule 6 (universality merge).} \emph{Two model nodes that share a calibration edge (the same observable fixes the same parameter of both) are merged into one node whose routes are the union of theirs.} Fires at Step~8, prepared at Step~7: the Fermi and muon-decay nodes share $G_F$, and the V$-$A node that replaces them is entered by the union of their routes plus the $\mu$-capture route. In the graph, parameter identity is node identity. Prediction: wherever two effective nodes are calibrated by one measured constant, a merge is due; the merged node has fewer parameters and more routes.

\paragraph{Rule 7 (absorption).} \emph{An effective node $m_{\rm eff}$ is absorbed by a fundamental node $m_{\rm fund}$ with $\dim[g]=0$ together with a new mediator sector $s$ such that (i) $m_{\rm eff}$ is the $q^2\ll M_s^2$ limit of $m_{\rm fund}$ (the matching condition is a lineage edge), (ii) the routes of $m_{\rm eff}$ become routes of $m_{\rm fund}$ with an extra leg $s(1)$, and (iii) the symmetry that makes $g$ dimensionless forces further model nodes whose observables carry badge~P.} Fires at Step~9: V$-$A $\to$ charged current, $G_F/\sqrt2=g^2/8M_W^2$, the four-leg routes become three-leg routes with $W(1)$, and the gauge group forces the neutral current and the self-couplings, with the neutral-current events, $M_W$, $M_Z$, $N_\nu$ and the triple-gauge couplings all predicted before observed. Absorption is the inverse of proliferation: the effective count in Figure~\ref{fig:s13} goes $0,1,2,1,0$ while the fundamental count rises monotonically.

\paragraph{Rule 8 (sector refinement).} \emph{The state layer changes only from outside itself: by Rule~1 (symmetry), by Rule~3 (a model node's back-edge), or by an observable that distinguishes sectors, in which case a sector is split, or replaced if the observable shows compositeness.} Fires at Steps~7 (the $\nu$ sector splits into $\nu_e$, $\nu_\mu$ after Danby et al.), 9 ($n,p$ replaced by $u,d$ after deep inelastic scattering) and 11 (the $\nu$ sectors must change because the massless representation cannot carry $\Delta m^2\neq0$). No sector was ever added by a model node alone; the W, Z and H sectors were added by absorption (Rule~7) and by the theory-internal case of Rule~2, and each was then confirmed by an observable with badge~P.

\paragraph{Rule 9 (calibration).} \emph{A model node with $p$ free parameters needs $p$ calibration edges (badge F); its predictive content is the number of its observable edges minus $p$.} In the final graph the calibration edges are $\alpha$ and $G_F$ (gauge--fermion cluster), $\sin^2\theta_W$ (fixed at Step~9), the fermion masses (Yukawa cluster) and $m_H$ (potential). Every other model$\to$observable edge is a prediction. The rule locates the predictive power of a theory in the graph: it is the excess of solid edges over dashed ones leaving a cluster.

\paragraph{Rule 10 (clustering).} \emph{When the model layer contains no effective node outside a fundamental one, the fundamental nodes are grouped by Lagrangian term; an absorbed effective node sits inside the cluster whose member reproduces it; routes and observable edges are attached to the cluster; and the number of clusters equals the number of terms of the Lagrangian.} Fires at Step~12. The clustered graph is the final form of a theory in GEN: it is what remains when the history is forgotten, and the lineage edges inside a cluster are the only trace of the history that the final form keeps.

\paragraph{The cycle.} Rules 2--7 compose into a loop that the graph runs twice. An observable with in-degree zero appears (Rule~2); a model node repairs it and, if needed, extends the state layer (Rule~3); the node's routes are read off its term (Rule~4); if the coupling is dimensionful the node is temporary (Rule~5); nodes that share a coupling merge (Rule~6); the merged node is absorbed by a fundamental one with a new sector, and the symmetry that makes the coupling dimensionless predicts further observables (Rule~7); the fundamental nodes cluster (Rule~10). The first run of the loop is anomaly-first (Steps~3--8, badges~A); its second half is prediction-first (Steps~9--10, badges~P); Step~11 opens a third run with a new in-degree-zero observable. Table~\ref{tab:rules} records where each rule fires, and the method view of Step~13 (Figures~\ref{fig:m4}--\ref{fig:m12}) shows the loop running: at each step the active phases of the cycle are filled and the operations applied to the graph are listed.

\begin{table}[htbp]\centering\small
\caption{Where each rule fires.}\label{tab:rules}
\begin{tabular}{@{}lll@{}}\toprule
Rule & Operation & Steps \\\midrule
1 symmetry partition & \textsc{partition} & 2 \\
2 in-degree-zero drive & \textsc{add} (model) & 3$\to$4, 5, 6, 7, 8, 11; theory-internal: 9, 10 \\
3 Fock extension & \textsc{extend} & 4 \\
4 route composition & routes & 5, 6, 7, 8, 9, 10, 11, 12 \\
5 effective status & marking & 6, 7, 8 \\
6 universality merge & \textsc{merge} & 8 (prepared in 7) \\
7 absorption & \textsc{absorb} + \textsc{add} sector & 9 \\
8 sector refinement & \textsc{split}, \textsc{replace} & 7, 9, 11 \\
9 calibration & F edges & 5, 6, 7, 9, 10 \\
10 clustering & \textsc{cluster} & 12 \\\bottomrule
\end{tabular}
\end{table}
